\chapter{Postgame and the Bonus Dungeon}
\label{written-postgame}

The ending opens another journey beneath the swamp outside
Evil Mountain. The Bonus Dungeon rewards exploration with
powerful equipment and the Drum of War, then leads to the
sleeping Estark in its deepest chamber.

Use the first visits to collect equipment and develop the
roster. A recovered party and carefully assigned support
items make the deeper return trips easier to sustain.
The optional confrontation provides a new test for the
family and monster companions raised through the adventure.

\section{Bonus Dungeon entrance and access prerequisites}

\subsection*{Loading the completed journey}

Finish Mildrath's two forms, the homeward town visits,
the Granvania celebration, and the credits through
\textbf{The End}. Reset after that screen and load
the save used for the completed game.

The ordinary adventure resumes from its pre-boss saved
state, with Bonus Dungeon access available. The final
treasures recorded before the ending remain useful
for this next journey.

\subsection*{The swamp entrance}

Cast \textbf{Return to Evil Mountain}. From the field
arrival point, walk \textbf{directly south into the swamp}
to enter the Bonus Dungeon.

The opening stair leads into a basement with several
branches. The first treasure pair lies west of the
arrival passage, while the onward northern route
eventually reaches a locked gate.

\subsection*{Preparation for the first treasure circuit}

Carry the \textbf{Final Key with an active member}.
Restore the intended fighters and reserves, and keep
Outside available for a recovery trip after collecting
the early rewards.

Assign the Sage's Stone and other battle-healing items
to members who can use them when needed. Keep stronger
healing, revival, and MP supplies distributed across
the travelling roster. The enemies are demanding enough
to make a fully prepared first entry worthwhile.

Save before descending, then explore with the aim of
retaining the new treasure and recording it during
the next town visit.

\begin{routecheck}[breakable=false]{Before entering the Bonus Dungeon}
\begin{description}
  \item[Completion:] The End reached and the associated save loaded.
  \item[Entrance:] Southern swamp outside Evil Mountain identified.
  \item[Key:] Final Key carried by an active traveller.
  \item[Party:] Active roles and reserves equipped, HP and MP
    restored, recovery items assigned, and departure saved.
\end{description}
\end{routecheck}

\section{Floor-by-floor route, branching passages, and puzzle solutions}

\subsection*{First basement: the entrance pair and Final Key gate}

From the arrival stairs, go west into the small side room.
Open both chests for a \textbf{Restoration Orb} and the
dungeon's fixed \textbf{Small Medal}.

\Achievement{61421}{Small Medal Collector}{25}{Collect 19 small medals from treasure chests}{
Update the chest-medal tally at this first western chest beside the
Restoration Orb. Check the unlock when the journey's total reaches
19 placed chest medals, including earlier collected chest finds.
Keep searched-location medals and monster or trapped-container battle
drops outside this tally. The threshold may have been reached earlier.}

Return to the arrival passage and continue north.

Pass the early eastern side path and follow the main
route to the northern fork. The northeastern branch
ends at the \textbf{first Mimic chest}. Prepare for
its optional battle if choosing to open it.

Return to the fork and take the northwestern route.
Continue west past the southern branch, then open
the \textbf{barred gate with the Final Key} and follow
the passage south into the next area.

\subsection*{The Shield of Radiance ledge}

The next area has large holes and narrow connecting
walkways. Work around the edges toward the southeastern
passage. A fall reaches a lower pocket, whose stairs
return to this area and provide another way to approach
the correct walkway.

Follow the passage east and south to the isolated chest
containing the \textbf{Shield of Radiance}.

Immediately after the pickup, equip the \textbf{Hero} with the
\textbf{Sun Crown}, \textbf{Dragon Staff}, \textbf{Royal Cloak},
and \textbf{Shield of Radiance}. Empty his accessory slot, removing
any Ring, armband, or Ribbon, and inspect all five equipment slots.

\Achievement{61424}{The Sacred Raiment}{5}{Equip the main character with the Sun Crown, Dragon Staff, Royal Cloak, and Shield of Radiance (and no accessories)}{
Equip the four retained pieces on the Hero as soon as this shield is
collected and leave his accessory slot empty. Check the unlock with
the complete set equipped.}

Retrace the ledge detour and continue southwest into the area
with cracked floor tiles.

\subsection*{Second basement treasure pocket: Drum and Sabre}

Use the stairs or deliberately fall through a cracked
tile to enter the lower treasure area. From the landing,
cross west through the swamp to the \textbf{Drum of War}
chest. Collect it and assign it to an available battle-item
user before continuing.

\AchievementSetup{61426}{That's Got To Hurt}{The Drum now provides
the attack boost for the single-hit procedure begun with Golem.
Develop Golem to level 11 for ChargeUp, equip a strong compatible
weapon, use the Drum or Bikill, then ChargeUp and attack one enemy.
Check one displayed hit for at least 500 damage. Keep ordinary
training, healing, and combat available during the attempt.}

Return toward the landing, then explore east and north
for the \textbf{Hell Sabre}, also called Hell Saber.
Climb back to the cracked-floor area.

Follow the intact route south, west, and north around
the holes to the northwestern stairs. Descend into
the first four-way maze.

\subsection*{Second basement maze: the cursed sword and onward stairs}

The arrival room has paths in four directions.
Take the eastern branch first and open the chest for
the \textbf{Sword of Decimation}. It is cursed.
Keep the ordinary weapon assigned while storing this find.

Return to the starting room and use the following
sequence of room exits:
\begin{enumerate}
  \item Go \textbf{west}, then \textbf{north}.
  \item Choose the \textbf{northwestern branch} and
    continue \textbf{north} again.
  \item Go \textbf{east}, then \textbf{south}.
  \item Go \textbf{east}, then \textbf{south} again.
  \item At the two-torch area, take the \textbf{eastern
    stairs downward}.
\end{enumerate}

The routes have deceptive screen connections. Keep the
starting room and the two-torch passage as landmarks
when returning after a detour.

\subsection*{Third basement maze: King Metal Armor}

The next maze's arrival room has a large central block.
Go \textbf{west, south, then south again} to the
\textbf{King Metal Armor} chest.

Retrace \textbf{north, north, then east} to the
central-block arrival room. Restore the party if the
detour's encounters have used substantial HP or MP.

\subsection*{The upper medicine pocket and three empty chests}

From the central-block room, go north and take the
upward stairs. They reach a separate second-basement
treasure pocket. One chest contains \textbf{Elven Medicine}
and the other is the \textbf{second Mimic}. Choose the
optional fight after preparing the active party.

Return down the stairs and south to the central-block
arrival room. Three eastern side rooms in this maze
hold \textbf{empty chests}: one east of the arrival
room, one east of the northern block room, and one
east of the southern block room.

For the onward route, leave the arrival room
\textbf{south, then east}. The chest in that eastern
room is empty, and its nearby stairs descend toward
the lower swamp passage.

\subsection*{Fourth basement: Iron Ball and the final descent}

Follow the long swamp corridor north, east, and south.
At the fork, take the western route. Explore north
for the \textbf{Iron Ball of Decimation} chest.

Return south one room and go west to the next stairs.
They lead into the fifth basement's lava chamber.
Follow the projecting walkway north to the sleeping
\textbf{Estark}.

Restore the intended combat group and assign its item
owners before speaking to him. The conversation provides
a choice to begin the optional battle, so the approach
is a convenient place to review the party first.

\begin{routecheck}{Bonus Dungeon route collections}
\begin{description}
  \item[Entrance side room:] Restoration Orb and one fixed Small Medal.
  \item[First-floor ledge:] Shield of Radiance.
  \item[Cracked-floor lower pocket:] Drum of War and Hell Sabre.
  \item[First four-way maze:] Unequipped cursed Sword of Decimation.
  \item[Second maze armor branch:] King Metal Armor.
  \item[Upper return pocket:] Elven Medicine and the separate Mimic.
  \item[Fourth basement spur:] Iron Ball of Decimation.
  \item[Other containers:] First Mimic at the northern entrance
    branch and three empty chests in the central-block maze.
\end{description}
\end{routecheck}

\section{Bonus treasure, the Drum of War, and unique equipment}

\subsection*{Assigning the new equipment}

Return to town after a successful treasure circuit,
recover, and save the new inventory. Inspect eligible
wearers and the equipment already in their slots
before assigning the rewards.

\begin{guidetable}[Bonus equipment and battle roles]
\begin{center}
\begin{tabular}{@{}p{0.24\linewidth}p{0.18\linewidth}p{0.45\linewidth}@{}}
\toprule
\GuideTableHeading{Reward} & \GuideTableHeading{Strength}
  & \GuideTableHeading{Assignment advice} \\
\midrule
Shield of Radiance & 65 defense & Fire and cold protection
  for a compatible shield wearer. \\
King Metal Armor & 95 defense & Strong body protection.
  Compare the current armor's elemental properties too. \\
Hell Sabre & 120 attack & A powerful focused weapon for
  an eligible physical member with a weaker weapon. \\
Iron Ball of Decimation & 125 attack & Whole-formation
  coverage for a compatible user during group encounters. \\
Sword of Decimation & Cursed & Retain as treasure and keep
  the ordinary battle weapon assigned. \\
\bottomrule
\end{tabular}
\end{center}
\end{guidetable}

The Hero's Dragon Staff already provides 125 attack,
and the Son has the Zenithian Sword. The new weapons
may give another active member or reserve a greater
improvement. Compare the actual compatible equipment
and the intended battle role.

\subsection*{The Drum's party-wide attack support}

Used through \textbf{Item during battle}, the
\textbf{Drum of War doubles the party's attack power}.
It can be used repeatedly and gives the physical
fighters a shared support action.

Assign it to a dependable active member and choose
the item action directly when starting a prepared
attack plan. A second member can apply Barrier or
handle healing while the Drum's user provides offense
support to the fighters.

When an opponent removes support effects, use the
Drum again during a safe recovery window. Treat the
Drum and Bikill as options for the same attack-boosting
job when distributing the turns, and keep essential
healing ahead of a repeated enhancement.

\subsection*{The recovery finds}

Keep the \textbf{Restoration Orb} and \textbf{Elven Medicine}
accessible in the party's inventory while assigning
recovery supplies. Carry the Orb with an intended battle
participant as a revival safeguard. The medicine provides
a substantial MP refill for a depleted caster. Inspect
the remaining supplies after each attempt before returning.

Continue carrying the Sage's Stone, Staff of Benediction,
revival supplies, and other MP recovery alongside the
new rewards. Their available users and remaining resources
determine how well the party sustains the next boss attempt.

\section{Rare companions and postgame recruitment routes}

\subsection*{The new Bonus Dungeon prospects}

The Bonus Dungeon provides recruitment encounters for
\textbf{Hellbattler}, \textbf{Lionex}, and \textbf{Springer}.
The repeated treasure and preparation journeys give the
party opportunities to meet them while keeping a recovery
stop within reach.

Hellbattler begins with \textbf{Explodet and Violent Flames}.
It learns \textbf{Bounce at level 7}, \textbf{Revive at 8},
\textbf{Blazemost at 11}, \textbf{Sparkling Breath at 13},
and \textbf{Infernal Flames at 15}. Its growing spell
and breath options let it cover several offensive and
support jobs.

Lionex has a \textbf{level-23 cap}. It begins with
Aerocross and develops Eerie Light at 8, Zap at 14,
and Lightning at 20. Inspect its current strength,
MP, and compatible equipment when choosing its position.

Springer has a \textbf{level-50 cap} and begins with
\textbf{Upper and Decrease}. Tornado at 7 and Aerocross
at 15 add offensive options to its initial defense-support
role.

\subsection*{Other encounters and the established roster}

Great Dragons, Attackbots, and Farewell Crags also appear
along the postgame routes. Their existing strengths can
complement the new equipment. Keep the strongest developed
members available while assessing a new recruit's reliability.

A King Healer with HealUsAll, Barrier, and Revive remains
valuable for the final encounter. Curer develops Revive
at 23 and Healusall at 30, providing another preparation
path for a recovery reserve.

Use manageable ordinary groups for training and recruitment,
then return to town to arrange the roster and save progress.
Choose the final combat party by its actual abilities,
protection, and recovery resources.

\section{Estark / Esterk: party setup, resistances, healing, and strategy}

\subsection*{The sleeping opponent and preparation}

\AchievementSetup{61413}{Underworld Challenge}{Prepare for a victory
in 25 turns or fewer, including a count of exactly 25. Use ordinary
training, strong equipment, Barrier, healing, and the Drum of War.
Save the prepared party and collected treasure at a normal church.
Hitoshikonomi, the early trap glitch, save-state cheating, and cheats
are excluded from this challenge. Check the spoken turn count and
achievement unlock after victory.}

At the end of the fifth-basement walkway, speak to
\textbf{Estark}, also called Esterk or Esturk.
He remembers his name and asks whether the party has
come to destroy him. Answer yes when ready to begin.

Estark has about \textbf{9,000 HP} and can act
\textbf{twice in a round}. A Hero and developed
combat companions around level 50 provide a useful
preparation target. Review the actual HP, MP, equipment,
and support roles before deciding that the group is ready.

Fire and cold protection are important against his
strong breath. Spell protection and a healthy HP
reserve help withstand Explodet. Status-protective
equipment such as an Elven Amulet is useful for a
member whose healing or item turn needs to remain available.

Use the first three positions for a durable attacker,
a dependable recovery user, and a flexible support or
second damage member. Equip the reserves for the same
pressure so they are ready to replace a threatened
active companion.

\begin{routecheck}{Estark preparation}
\begin{itemize}
  \item Restore the intended fighters and reserves before
    the approach, and save the bonus equipment already collected.
  \item Assign the Drum of War and Sage's Stone to available
    users, with the Staff of Benediction as another recovery option.
  \item Keep Barrier available through the Son, King Healer,
    or another developed support member.
  \item Bring strong group healing and reliable revival,
    alongside ordinary item recovery.
  \item Distribute Magic Water, Rings of Prayer, Elven Medicine,
    and revival supplies across the travelling roster.
  \item Check fire, cold, spell, and status protection on
    every member expected to enter the battle.
\end{itemize}
\end{routecheck}

\subsection*{The opening support turns}

Open with \textbf{Barrier} and the \textbf{Drum of War}
when the party's condition allows both support actions.
The remaining member can attack or handle an immediate
injury. Choose the Drum's item command directly.

Bring the supported physical fighters forward as the
lineup settles. A developed Golem, another strong monster,
the Hero, or the Son can provide sustained attacks,
according to the equipment and abilities currently available.

Keep one dependable recovery action in the round plan.
The Sage's Stone can handle routine group damage, and
the staff helps single-target recovery while preserving
the caster's MP for stronger spells or revival.

\subsection*{Two-action pressure and the healing reserve}

Estark can combine \textbf{Explodet with powerful breath}
or another dangerous attack in the same round. Barrier
reduces the breath portion, while the spell damage still
requires sufficient HP and equipment protection.

Restore endangered companions before returning to an
all-out attack turn. Keep HP well above 200 when the
group's maximum HP allows it, and judge the margin
against the damage the current equipment is taking.

Rotate King Healer or another strong recovery member
in for a group refill when needed. An injured active
fighter can move to reserve while a prepared replacement
continues the offensive role. Bring a revived companion
back into the lineup after its condition is suitable.

His third-eye action can put companions to sleep.
Keep a functioning recovery user available through
the reserves when an important active member loses
its turn. The planned status protection supports that
recovery continuity.

\subsection*{Restoring enhancements and handling reflection}

Estark can clear the party's magical enhancements.
Renew \textbf{Barrier} when breath pressure is the
immediate threat, then use the \textbf{Drum} again
on a safe turn to restore the attack plan.

Increase can help when physical blows are causing
trouble. Give urgent healing priority over rebuilding
an enhancement while a member is close to falling.

Watch for Estark's \textbf{Bounce}. Use physical
attacks against a reflected target and account for
reflection before choosing an offensive spell.
The carried recovery items remain useful to their
assigned active users.

Use MP supplies before the essential caster's reserve
becomes too low for protection, group healing, or revival.
Another member can take that item turn while the regular
healer keeps supporting the party.

\subsection*{Victory and the spoken turn count}

Maintain the protected attack and recovery pattern
until Estark is defeated. Victory awards
\textbf{5,000 experience} and \textbf{510 G}, with
a \textbf{Small Medal} possible as a battle reward.

Estark announces how many turns the victory took,
then disappears. Record the number if using the
encounter as a personal performance goal. Return
to town, recover, and save the retained treasure
and training progress.

\Achievement{61412}{Until Next Time}{25}{Defeat Esturk}{
Defeat Estark and finish his victory dialogue. Check the unlock after
the battle and retain the party's ordinary saved progress.}

\Achievement{61413}{Underworld Challenge}{25}{Defeat Esturk in 25 turns or less (no cheating)}{
After victory, check Estark's announced turn count for 25 or fewer.
A count of exactly 25 qualifies. Confirm the unlock from the
legitimate prepared attempt before another run. Use ordinary training,
the Drum, Barrier, and equipment throughout, with Hitoshikonomi,
the early trap glitch, save-state cheating, and cheats excluded.}

\section{Repeat battles, challenge goals, and remaining completion tasks}

\subsection*{Returning for another attempt}

Leave the dungeon and return to Estark's chamber
to face him again. The spoken victory count provides
a repeatable score for improving the party's preparation
and turn choices.

For \textbf{Underworld Challenge}, aim for \textbf{25 turns or fewer}
on a legitimate repeat victory and confirm the RetroAchievements
unlock after the spoken count. Exactly 25 turns meets the goal.

Begin with a reliable clear, then refine one part
of the plan at a time. Compare the opening support
turns, the Drum user's assignment, reserve transitions,
MP refills, and how quickly protection is renewed
after a clearing action.

Track both the successful turn count and the recovery
resources used. A faster attempt can remain stable
when its fighters, healing, and support are balanced.

\subsection*{Finishing the retained collection}

Review the nine fixed rewards and the party's storage
before planning another treasure trip. Keep optional
Mimic rewards separate from the single placed Small Medal.
The three empty maze chests complete the container search.

Use the new equipment on the compatible companions
who benefit most, and save after a meaningful upgrade
or recruitment result. The Drum supports ordinary
training journeys as well as the repeat boss fight.

The family and monsters can also return to Mildrath
with the postgame equipment for another ending tour.
Restore the roster first and save the new treasure
before the two-form sequence.

\begin{routecheck}[breakable=false]{Postgame completion}
\begin{description}
  \item[Fixed rewards:] Restoration Orb, Small Medal,
    Shield of Radiance, Drum of War, Hell Sabre,
    Sword of Decimation, King Metal Armor,
    Elven Medicine, and Iron Ball of Decimation.
  \item[Containers:] Both Mimic locations and all three
    empty maze chests accounted for.
  \item[Assignments:] New compatible equipment and the
    Drum's active user selected, with cursed gear retained safely.
  \item[Party:] Intended recruitment and training progress
    saved, reserves equipped, and recovery roles prepared.
  \item[Boss:] Estark defeated and the victory turn count
    recorded for any future challenge attempt.
  \item[Save:] Treasure and preparation retained before
    another Estark battle or ending run.
\end{description}
\end{routecheck}

